\en{In this chapter, which can be read straight after chapter~\ref{\en{EN}kapGitter}, we prove the Picard Fuchs differential equation.
The proof follows \cite[p.~33-34, ch.~I.2, §3]{\en{EN}klein} -- but there every $g$ has a different sign.}%
\de{In diesem Kapitel, das direkt im Anschluss an Kapitel~\ref{\en{EN}kapGitter} gelesen werden kann, beweisen wir die Picard-Fuchs-Differentialgleichung.
Der Beweis orientiert sich stark an \cite[S.~33-34, Kap.~I.2, §3]{\en{EN}klein}. Dort hat allerdings jedes $g$ genau das andere Vorzeichen.}\\


\begin{theo}[Picard\en{ }\de{-}Fuchs]\label{\en{EN}picardfuchs}
\en{The periods $\Omega_{1,2}(J)$ of $L_J$ from Def.~\ref{\en{EN}definLJ} are solutions to the following differential equation:}%
\de{Für die Perioden $\Omega_{1,2}(J)$ von $L_J$ aus Def.~\ref{\en{EN}definLJ} gilt folgende Differentialgleichung:}
$$\frac{d^2\Omega}{dJ^2}+\frac 1 J \cdot \frac{d\Omega}{dJ} + \frac{31J-4}{144J^2(J-1)^2}\cdot\Omega = 0$$
\end{theo}
\begin{proof} 
\en{We start with the representations of the basic periods $\Omega_{1,2}$ and basic quasiperiods $H_{1,2}$ of the lattice $L_J$ from Prop.~\ref{\en{EN}satzint}:}%
\de{Zunächst gilt nach Satz~\ref{\en{EN}satzint} für die Perioden $\Omega_{1,2}$ und Quasiperioden $H_{1,2}$ des Gitters $L_J$:}
\begin{align*}
    \Omega_k &= \oint_{\alpha_k} \frac{dx}{y}\qquad\text{\en{and}\de{und}}\qquad H_k = -\oint_{\alpha_k} \frac{x~dx}{y}
\end{align*}
\en{The defining equation of $X(L_J)$ reads $y^2=4x^3-g(x+1)$ with $g=\frac{27J}{J-1}$ (cf.~Prop.~\ref{\en{EN}satz44}). This yields:}%
\de{Die definierende Gleichung von $X(L_J)$ lautet $y^2=4x^3-g(x+1)$ mit $g=\frac{27J}{J-1}$ (vgl.~Satz~\ref{\en{EN}satz44}). Hieraus folgt:}
\begin{align}
    \frac{d}{dg}(y^2) &= \frac{d}{dg}\lk\left(\frac 1 y\right)^{-2}\rk
    = -2 \cdot \left(\frac 1 y\right)^{-3} \cdot \frac{d}{dg}\lk\frac 1 y\rk\label{\en{EN}ablxyg}\\
    \Longrightarrow\quad\frac{d}{dg}\lk\frac 1 y\rk &= \frac{-1}{2y^3}\cdot \frac{d}{dg}(y^2) 
    = \frac{-1}{2y^3}\cdot \frac{d}{dg}(4x^3-g(x+1)) = \frac{x+1}{2y^3}\nonumber
\end{align}
\en{From this we obtain the derivatives of $\Omega$ and $H$ with respect to $g$ using Leibniz's rule:}%
\de{Somit folgen die Ableitungen von $\Omega$ und $H$ nach $g$ mit der Leibnizregel:}
\begin{align}
    \frac{d\Omega}{dg} &= \frac{d}{dg}\lk\oint_{\alpha} \frac{dx}{y}\rk
                        = \oint_{\alpha}\frac{d}{dg}\lk\frac{1}{y}\rk dx
                        = \oint_{\alpha} \frac{x~dx}{2y^3} + \oint_{\alpha} \frac{dx}{2y^3}\label{\en{EN}stern}\\
    \text{\en{and}\de{und} }\quad\frac{dH}{dg} &= \frac{d}{dg}\lk-\oint_{\alpha} \frac{x~dx}{y}\rk
                        = -\oint_{\alpha}\frac{d}{dg}\lk\frac{1}{y}\rk\cdot x~dx
                        = -\oint_{\alpha} \frac{x^2~dx}{2y^3} - \oint_{\alpha} \frac{x~dx}{2y^3}\nonumber
\end{align}
\en{We have to calculate the values $I_n:=\oint_{\alpha} \frac{x^n~dx}{2y^3}$ for $n=0,1,2$.
In order to do this, we use the functions $f_n(x):=\frac{x^n}{y}$ with $n=0,1,2$ and start by deriving $f_0(x)=\frac 1 y$ as in~(\ref{\en{EN}ablxyg}):}%
\de{Wir müssen also noch die Werte der Integrale $I_n:=\oint_{\alpha} \frac{x^n~dx}{2y^3}$ für $n=0,1,2$ berechnen.
Hierfür verwenden wir die Funktionen $f_n(x):=\frac{x^n}{y}$ mit $n=0,1,2$, wobei wir zunächst $f_0(x)=\frac 1 y$ so wie in~(\ref{\en{EN}ablxyg}) ableiten:}
\begin{align*}
    f_0 (x) &=\frac 1 y\\
    f_0'(x) &= \frac{d}{dx}\lk\frac 1 y\rk = \frac{-1}{2y^3}\cdot\frac{d}{dx}(y^2) = \frac{-1}{2y^3}\cdot \frac{d}{dx}(4x^3-g(x+1))\\
    &= \frac{-1}{2y^3}\cdot(12x^2-g)=\frac{g-12x^2}{2y^3}\\
    f_1 (x) &=\frac x y = x \cdot f_0(x)\\
    f_1'(x) &=f_0(x)+x\cdot f_0'(x) = \frac 1 y + x \cdot \frac{g-12x^2}{2y^3} = \frac 1 y + \frac{gx-3\cdot 4x^3}{2y^3}\\
            &= \frac 1 y + \frac{gx-3\cdot \left(y^2 + g x +g\right)}{2y^3}
            = \frac 1 y + \frac{-3y^2 - 2 g x -3 g}{2y^3}\\ &= \frac 2 {2y} - \frac{3y^2}{2y^3} - \frac{2gx+3g}{2y^3}
            = -\frac {1} {2y} - \frac{2gx+3g}{2y^3}\end{align*}\begin{align*}
    f_2 (x) &= \frac{x^2}{y} = x\cdot f_1(x)\\
    f_2'(x) &= f_1(x) + x \cdot f_1'(x) = \frac x y + x \cdot \left(-\frac {1} {2y} - \frac{2gx+3g}{2y^3}\right)\\
            &=\frac {2x} {2y} -\frac {x} {2y} - \frac{2gx^2+3gx}{2y^3} = \frac {x} {2y} - \frac{2gx^2+3gx}{2y^3}
\end{align*}

\en{Prop.~\ref{\en{EN}satzint} tells us that $\alpha$ is a closed path, which avoids the zeros and poles of $\wp$ (cf.~Remark~\ref{\en{EN}bemwege}). From this we deduce that}%
\de{Aus Satz~\ref{\en{EN}satzint} folgt, dass $\alpha$ eine geschlossene Kurve ist, die die Nullstellen und Polstellen der $\wp$-Funktion vermeidet (vgl.~Bem.~\ref{\en{EN}bemwege}). Hieraus folgt:}
$$\oint_{\alpha} f_n'(x)dx=\int_0^1f_n'(\alpha(t))\alpha'(t)dt=\left[f_n(\alpha(t))\right]_0^1=0$$
\en{Using the derivatives of $f_0$, $f_1$ and $f_2$ we calculated above we get the following relations between the values $I_n=\oint_{\alpha} \frac{x^n~dx}{2y^3}$:}%
\de{Wenn wir die oben berechneten Ableitungen von $f_0$, $f_1$ und $f_2$ verwenden erhalten wir die folgenden Beziehungen zwischen den Integralen $I_n=\oint_{\alpha} \frac{x^n~dx}{2y^3}$:}
\begin{align*}
    0= \oint_\alpha f_0'(x)dx &= \qquad\qquad\qquad~ g \cdot \oint_\alpha \frac {dx} {2y^3}\quad ~~- 12 \cdot \oint_\alpha \frac {x^2~dx} {2y^3}\\
    0= \oint_\alpha f_1'(x)dx &= - \underbrace{\oint_\alpha \frac {dx} {2y}}_{=~\frac 1 2 \Omega}\quad  - ~ 2g \cdot \oint_\alpha \frac {x~dx} {2y^3}\quad - 3g\cdot \oint_\alpha \frac {dx} {2y^3}\\
    0=\oint_\alpha f_2'(x)dx &= \underbrace{\oint_\alpha \frac {x~dx} {2y}}_{=~-\frac 1 2 H}\quad - ~2g \cdot \oint_\alpha \frac {x^2~dx} {2y^3}\quad  -3g \cdot \oint_\alpha \frac {x~dx} {2y^3}
\end{align*}
\en{This produces the following system of linear equations in $I_0,I_1,I_2$:}%
\de{Also gilt für die Integrale $I_0,I_1,I_2$ folgendes lineares Gleichungssystem:}
\begin{align*}
    \left|
    \begin{aligned}
        \phantom{3}g\cdot I_0 \phantom{~+2g\cdot I_1}-12\cdot I_2 &= 0 & \text{(I)}~\\
        ~3 g \cdot I_0 + 2g\cdot I_1 \phantom{-12g\cdot I_2} &= -\frac 1 2 \cdot \Omega & \text{(II)}~\\
        \phantom{3g\cdot I_0 +} 3g\cdot I_1 + 2g\cdot I_2 &= -\frac 1 2\cdot H & \text{(III)}~
    \end{aligned}
    \right|
    \quad \Longrightarrow \quad 
    \left|
    \begin{aligned}
        I_0 &= \displaystyle\frac{9\Omega-6H}{2g(g-27)}\\[1ex]
        ~I_1 &= \displaystyle \frac{18H-g\Omega}{4g(g-27)}~\\[1ex]
        I_2 &= \displaystyle\frac{3\Omega-2H}{8(g-27)}
    \end{aligned}
    \right|
\end{align*}
\en{The value of $I_2$ comes from (III) -- 1\ko5 $\cdot$ (II) + 4\ko5 $\cdot$ (I). Then, (I) yields $I_0$ and (II) yields $I_1$.
Now we use these results in~(\ref{\en{EN}stern}) to get the desired values of $\frac{d\Omega}{dg}$ and $\frac{dH}{dg}$:}%
\de{Den Wert von $I_2$ erhält man z.B.~aus (III) -- 1\ko5 $\cdot$ (II) + 4\ko5 $\cdot$ (I). Dann erhält man aus $I_0$ aus (I) und dann $I_1$ aus (II).
Das setzen wir in~(\ref{\en{EN}stern}) ein und erhalten $\frac{d\Omega}{dg}$ und $\frac{dH}{dg}$:}
\begin{align*}
    \frac{d\Omega}{dg} &= I_0 + I_1 = \frac{9\Omega-6H}{2g(g-27)}+\frac{18H-g\Omega}{4g(g-27)} = \frac{(18-g)\Omega +6H}{4g(g-27)}\\
    \text{\en{and}\de{und} }\quad\frac{dH}{dg} &= -I_1 - I_2 = -\frac{18H-g\Omega}{4g(g-27)}-\frac{3\Omega-2H}{8(g-27)} = \frac{(2g-36)H-g\Omega}{8g(g-27)}
\end{align*}
\en{Using $g=\frac{27J}{J-1}$ we can transform these equations in $g$ into equations in $J$.
For this transformation, we use $\frac{dg}{dJ} = \frac{-27}{(J-1)^2}$ and $\frac{d}{dg} = \left(\frac{dg}{dJ}\right)^{-1}\cdot\frac{d}{dJ}$.
This transforms the two equations into}%
\de{Mit $g=\frac{27J}{J-1}$ kann man diese beiden Gleichungen, die von $g$ abhängen, in Gleichungen umwandeln, die stattdessen von $J$ abhängen.
Dazu nutzen wir $\frac{dg}{dJ} = \frac{-27}{(J-1)^2}$ und dass wir die Ableitung nach $g$ somit durch eine nach $J$ ersetzen können: $\frac{d}{dg} = \left(\frac{dg}{dJ}\right)^{-1}\cdot\frac{d}{dJ}$

Deshalb übersetzen sich die beiden Gleichungen in}
\begin{align*}
    \frac{(J-1)^2}{-27}\cdot\frac{d\Omega}{dJ}
    &= \frac{\left(18-\frac{27J}{J-1}\right)\Omega +6H}{4\cdot\frac{27J}{J-1}\cdot\left(\frac{27J}{J-1}-27\right)}\\
    \text{\en{and}\de{und}}\qquad\frac{(J-1)^2}{-27}\cdot\frac{dH}{dJ}
    &= \frac{\left(2\cdot\frac{27J}{J-1}-36\right)H-\frac{27J}{J-1}\Omega}{8\cdot\frac{27J}{J-1}\cdot\left(\frac{27J}{J-1}-27\right)}
\end{align*}
\en{which simplifies to:}\de{was sich wie folgt vereinfachen lässt:}
\begin{align}
    36 J(J-1)\frac{d\Omega}{dJ} &= 3(J+2)\Omega -2(J-1)H\label{\en{EN}dOdJ}\\
    \intertext{\en{and}\de{und}} 24 J(J-1)\frac{dH}{dJ} &= 3J \Omega -2(J+2)H\label{\en{EN}dHdJ}
\end{align}
\en{Now we derive equation~(\ref{\en{EN}dOdJ}) by $J$ once more (using the product rule), and obtain after grouping similar terms:}%
\de{Wenn wir nun Gleichung~(\ref{\en{EN}dOdJ}) nochmals nach $J$ ableiten (unter Beachtung der Produktregel), erhalten wir nach Zusammenfassen gleichartiger Terme:}
\begin{align*}
    36J(J-1)\frac{d^2\Omega}{dJ^2}+(69J-42)\frac{d\Omega}{dJ}+2(J-1)\frac{dH}{dJ}-3\Omega+2H&=0
\end{align*}
\en{Next we multiply this equation with $12J$ and eliminate $\frac{dH}{dJ}$ using~(\ref{\en{EN}dHdJ}). This yields:}%
\de{Nun multiplizieren wir diese Gleichung mit $12J$ und eliminieren $\frac{dH}{dJ}$ mit Hilfe von~(\ref{\en{EN}dHdJ}). Das führt auf:}
\begin{align*}
    432 J^2(J-1)\cdot\Omega'' + 12 J\cdot(69J-42)\cdot\Omega' - 33 J \cdot \Omega + (11J-2) \cdot 2H = 0
\end{align*}
\en{Here, we multiply with $(J-1)$ and eliminate $H$ using~(\ref{\en{EN}dOdJ}). This yields:}%
\de{Das multiplizieren wir mit $(J-1)$ und eliminieren $H$ mit Hilfe von~(\ref{\en{EN}dOdJ}). Wir erhalten:}
\begin{align*}
    432 J^2(J-1)^2\cdot\Omega'' + 12 J(J-1)(36J-36)\cdot\Omega' + (93J-12) \cdot \Omega = 0
\end{align*}
\en{One last division by $432 J^2(J-1)$ yields the Picard Fuchs differential equation:}%
\de{Eine letzte Division durch $432 J^2(J-1)^2$ liefert die Picard-Fuchs-Differentialgleichung:}\belowdisplayskip=-12pt
\begin{align*}
    \frac{d^2\Omega}{dJ^2} + \frac 1 J \cdot \frac{d\Omega}{dJ} + \frac{31J-4}{144J^2(J-1)^2}\cdot\Omega &= 0
\end{align*}
\end{proof}
